%% toc: P1 6 6.7
\section{Proofs for soundness, completeness and consistency}
\label{sec:annex-p01-w02-p01-019-consistency}

Notation, standing data and definitions are those of
\cref{sec:p01-w02-p01-019-consistency}. For a set \(\Gamma\) and a datum
\(d\), \(d\entails_{\mathcal S}\Gamma\) means that \(d\) satisfies every
member of \(\Gamma\). Every datum satisfies the empty set.

\begin{annexitem}{annex:p01-w02-p01-019-forms}{thm:p01-w02-p01-019-forms}{Comparison of the forms of consistency}
\begin{proof}
\emph{Clause 0.} Suppose \(\Gamma\subseteq\Delta\) and
\(\Gamma\derives_{\mathcal C}\varphi\). By the definition of derivability in
\cref{sec:levels-03} there is a derivation in \(\mathcal C\) with conclusion
\(\varphi\) whose undischarged assumptions are members of \(\Gamma\). They are
then members of \(\Delta\), so the same derivation witnesses
\(\Delta\derives_{\mathcal C}\varphi\). Let \(\Gamma\subseteq\Delta\). If
\(\Gamma\) is inconsistent, some \(\varphi\) satisfies
\(\Gamma\derives_{\mathcal C}\varphi\) and
\(\Gamma\derives_{\mathcal C}\neg\varphi\), and by monotonicity \(\Delta\)
derives both, so \(\Delta\) is inconsistent. If \(\Gamma\) is trivial,
monotonicity gives \(\Delta\derives_{\mathcal C}\psi\) for every \(\psi\), so
\(\Delta\) is trivial. If \(\Gamma\derives_{\mathcal C}\bottom\), then
\(\Delta\derives_{\mathcal C}\bottom\). Contraposing each statement, every
subset of a consistent (absolutely consistent, \(\bottom\)-consistent) set is
consistent (absolutely consistent, \(\bottom\)-consistent).

\emph{Clause 1.} Suppose \(\Gamma\) is not absolutely consistent, so it
derives every formula. Since \(F\) is nonempty, choose \(\varphi\in F\). Since
\(F\) is closed under \(\neg\), \(\neg\varphi\in F\). Then \(\Gamma\) derives
both \(\varphi\) and \(\neg\varphi\) and is inconsistent. Contraposing, every
consistent set is absolutely consistent.

\emph{Clause 2.} Suppose \(\mathcal C\) has explosion and \(\Gamma\) is
inconsistent, witnessed by \(\varphi\). For every \(\psi\), explosion applied
to \(\Gamma\derives_{\mathcal C}\varphi\) and
\(\Gamma\derives_{\mathcal C}\neg\varphi\) gives
\(\Gamma\derives_{\mathcal C}\psi\), so \(\Gamma\) is trivial. Contraposing,
every absolutely consistent set is consistent. With clause 1, the two forms
coincide.

\emph{Clause 3.} Suppose \(\mathcal C\) has reductio, and let
\(\Gamma\derives_{\mathcal C}\varphi\) and
\(\Gamma\derives_{\mathcal C}\neg\varphi\). Let \(\psi\in F\). By clause 0,
\(\Gamma\cup\{\neg\psi\}\) derives \(\varphi\) and \(\neg\varphi\). Reductio,
applied with the formula \(\psi\) and the witness \(\varphi\), gives
\(\Gamma\derives_{\mathcal C}\psi\). Hence \(\mathcal C\) has explosion.

\emph{Clause 4.} Suppose \(\bottom\in F\) and (B1), (B2) hold. If
\(\Gamma\derives_{\mathcal C}\varphi\) and
\(\Gamma\derives_{\mathcal C}\neg\varphi\), then (B2) gives
\(\Gamma\derives_{\mathcal C}\bottom\), and (B1) gives
\(\Gamma\derives_{\mathcal C}\psi\) for every \(\psi\). So \(\mathcal C\) has
explosion, and by clause 2 absolute consistency and consistency coincide. It
remains to compare \(\bottom\)-consistency with absolute consistency. If
\(\Gamma\derives_{\mathcal C}\bottom\), (B1) makes \(\Gamma\) trivial.
Conversely, if \(\Gamma\) is trivial, it derives \(\bottom\in F\). Hence
\(\Gamma\) is \(\bottom\)-inconsistent iff it is trivial, which is the
claimed equivalence.
\end{proof}
\end{annexitem}

\begin{annexitem}{annex:p01-w02-p01-019-ndcpl}{thm:p01-w02-p01-019-ndcpl}{The structural properties of \(\mathsf{ND}_{\mathsf{CPL}}\)}
\begin{proof}
Derivations are as in the definition of a natural-deduction derivation in
\cref{sec:prop-deduction-01}. They are finite rooted trees
with marked assumption leaves, discharge labels from \(\N\), injective label
renaming, and the merging relabeling.

\emph{Reflexivity.} For \(\varphi\in\Gamma\), the tree with one node, an
assumption leaf \(\varphi\), has conclusion \(\varphi\). Its only open
assumption is \(\varphi\in\Gamma\).

\emph{(B1).} Let \(D\) derive \(\bottom\) with open assumptions in
\(\Gamma\). Adding a root \(\psi\) below the root of \(D\) by \(\bottom E\)
gives a derivation of \(\psi\). The rule \(\bottom E\) discharges nothing, so
its open assumptions are those of \(D\), all in \(\Gamma\).

\emph{(B2).} Let \(D_1\) derive \(\varphi\) and \(D_2\) derive
\(\neg\varphi\), each with open assumptions in \(\Gamma\). Rename the labels
of \(D_2\) injectively to a set disjoint from the labels of \(D_1\). This
renaming preserves \(D_2\) and the occurrences it discharges. Join the two
roots by \(\neg E\) to conclude \(\bottom\). The rule \(\neg E\) discharges
nothing. Because the label sets are disjoint, no discharge in either
subderivation reaches into the other. So the open assumptions of the new
derivation are those of \(D_1\) and \(D_2\), all in \(\Gamma\).

\emph{Reductio.} Suppose \(\Gamma\cup\{\neg\varphi\}\) derives some \(\psi\)
and \(\neg\psi\). By (B2) there is a derivation \(D\) of \(\bottom\) whose
open assumptions lie in \(\Gamma\cup\{\neg\varphi\}\). There are two cases.
\begin{itemize}
  \item No open assumption occurrence of \(D\) is an occurrence of
        \(\neg\varphi\). Then every open assumption of \(D\) lies in
        \(\Gamma\), so \(\Gamma\derives\bottom\), and (B1) gives
        \(\Gamma\derives\varphi\).
  \item Some open occurrence of \(\neg\varphi\) remains. Mark all open
        occurrences of \(\neg\varphi\) in \(D\) with one label \(i\) fresh
        from every label of \(D\). By the merging relabeling this preserves
        the derivation and the discharges of each of its rule applications.
        Apply \(\mathsf{RAA}\) with label \(i\) at the root. Its premise is
        \(\bottom\), its conclusion is \(\varphi\), and it discharges exactly
        the occurrences marked \(i\), of which there is at least one, as the
        nonvacuous-discharge requirement demands. Every remaining open
        occurrence is an open occurrence of \(D\) of a formula other than
        \(\neg\varphi\), hence a member of \(\Gamma\). So
        \(\Gamma\derives\varphi\).
\end{itemize}
The final statement now follows from clause 4 of
\cref{thm:p01-w02-p01-019-forms}, since \(\bottom\in\mathsf{Fm}_{\mathsf{CPL}}\)
and the falsum properties hold.
\end{proof}
\end{annexitem}

\begin{annexitem}{annex:p01-w02-p01-019-soundness}{thm:p01-w02-p01-019-soundness}{Soundness and consistency}
\begin{proof}
\emph{Clause 1.} Let \(d\entails_{\mathcal S}\Gamma\), and suppose \(\Gamma\)
is inconsistent, with \(\Gamma\derives_{\mathcal C}\varphi\) and
\(\Gamma\derives_{\mathcal C}\neg\varphi\). Soundness gives
\(\Gamma\entails_{\mathcal S}\varphi\) and
\(\Gamma\entails_{\mathcal S}\neg\varphi\). Since \(d\) satisfies \(\Gamma\),
it satisfies \(\varphi\) and \(\neg\varphi\), contrary to exclusive negation.
So \(\Gamma\) is consistent, and by clause 1 of
\cref{thm:p01-w02-p01-019-forms} absolutely consistent. If \(\bottom\in F\),
no datum satisfies \(\bottom\), and \(\Gamma\derives_{\mathcal C}\bottom\),
then soundness gives \(\Gamma\entails_{\mathcal S}\bottom\). The datum \(d\)
would satisfy \(\bottom\), which is impossible. So
\(\Gamma\notderives_{\mathcal C}\bottom\).

\emph{Clause 2.} Choose \(d\in\mathcal D_{\mathcal S}\). It satisfies
\(\varnothing\), so \(\varnothing\) is satisfiable, and by clause 1 it is
consistent. That is the consistency of \(\mathcal C\).

\emph{Clause 3.} It is the contrapositive of the first sentence of clause 1.

\emph{Converse.} Assume the classical negation clause, reflexivity, and that
every satisfiable set is consistent. Let \(\Gamma\derives_{\mathcal C}\varphi\).
By clause 0 of \cref{thm:p01-w02-p01-019-forms},
\(\Gamma\cup\{\neg\varphi\}\derives_{\mathcal C}\varphi\). By reflexivity,
\(\Gamma\cup\{\neg\varphi\}\derives_{\mathcal C}\neg\varphi\). So
\(\Gamma\cup\{\neg\varphi\}\) is inconsistent, and by the hypothesis it is
unsatisfiable. Clause 1 of \cref{thm:p01-w02-p01-016-unsat-reduction}, which
uses the classical negation clause, gives \(\Gamma\entails_{\mathcal S}\varphi\).
Hence \(\mathcal C\) is sound. The classical clause implies exclusive
negation, so clause 1 supplies the forward direction of the stated
equivalence.
\end{proof}
\end{annexitem}

\begin{annexitem}{annex:p01-w02-p01-019-completeness}{thm:p01-w02-p01-019-completeness}{Completeness and consistency}
\begin{proof}
\emph{Clause 1.} Let \(\mathcal C\) be complete for \(\mathcal S\), and
suppose \(\Gamma\) is unsatisfiable. By clause 4 of
\cref{thm:p01-w02-p01-016-unsat-reduction}, which holds for every semantics,
\(\Gamma\entails_{\mathcal S}\psi\) for every \(\psi\in F\). Completeness
gives \(\Gamma\derives_{\mathcal C}\psi\) for every \(\psi\), so \(\Gamma\) is
trivial. Contraposing, every absolutely consistent set is satisfiable. Every
consistent set is absolutely consistent by clause 1 of
\cref{thm:p01-w02-p01-019-forms}, and so it is satisfiable.

\emph{Clause 2.} Assume the classical negation clause, reductio, and that
every consistent set is satisfiable. Let \(\Gamma\entails_{\mathcal S}\varphi\).
By clause 1 of \cref{thm:p01-w02-p01-016-unsat-reduction},
\(\Gamma\cup\{\neg\varphi\}\) is unsatisfiable. By the hypothesis it is
therefore not consistent, so some \(\psi\) has
\(\Gamma\cup\{\neg\varphi\}\derives_{\mathcal C}\psi\) and
\(\Gamma\cup\{\neg\varphi\}\derives_{\mathcal C}\neg\psi\). Reductio gives
\(\Gamma\derives_{\mathcal C}\varphi\). Hence \(\mathcal C\) is complete.

\emph{Clause 3.} Assume the classical negation clause and reductio. Clause 1
shows that completeness implies that every absolutely consistent set is
satisfiable. By clause 3 of \cref{thm:p01-w02-p01-019-forms}, \(\mathcal C\)
has explosion, so by clause 2 of that lemma the consistent sets are the
absolutely consistent sets. Hence the second and third conditions are the
same condition. Clause 2 shows that the second condition implies
completeness.
\end{proof}
\end{annexitem}

\begin{annexitem}{annex:p01-w02-p01-019-adequate}{thm:p01-w02-p01-019-adequate}{Consistency and deductive completeness for an adequate calculus}
\begin{proof}
Since \(\mathcal C\) is sound and complete, \(\Gamma\derives_{\mathcal C}\varphi\)
iff \(\Gamma\entails_{\mathcal S}\varphi\) for all \(\Gamma,\varphi\).

\emph{Clause 1.} If \(\Gamma\) is satisfiable, it is consistent by clause 1 of
\cref{thm:p01-w02-p01-019-soundness}, since the classical clause gives
exclusive negation. If \(\Gamma\) is consistent, it is absolutely consistent
by clause 1 of \cref{thm:p01-w02-p01-019-forms}. If \(\Gamma\) is absolutely
consistent, it is satisfiable by clause 1 of
\cref{thm:p01-w02-p01-019-completeness}. The three conditions form a cycle of
implications and are therefore equivalent.

\emph{Clause 2.} By clause 1, \(\mathcal C\) is consistent iff
\(\varnothing\) is satisfiable. Every datum satisfies \(\varnothing\), so
\(\varnothing\) is satisfiable iff
\(\mathcal D_{\mathcal S}\neq\varnothing\).

\emph{Clause 3.} By clause 1, consistency of \(\Gamma\) is satisfiability. By
the displayed equivalence, \(\Gamma\) is deductively complete iff, for every
\(\varphi\), \(\Gamma\entails_{\mathcal S}\varphi\) or
\(\Gamma\entails_{\mathcal S}\neg\varphi\).

(\(\Rightarrow\)) Let \(\Gamma\) be satisfiable and deductively complete. Let
\(d,d'\) satisfy \(\Gamma\) and let \(\varphi\in F\). If
\(\Gamma\entails_{\mathcal S}\varphi\), both satisfy \(\varphi\). Otherwise
\(\Gamma\entails_{\mathcal S}\neg\varphi\). Then both satisfy
\(\neg\varphi\), and by the classical clause neither satisfies \(\varphi\).
In both cases \(d\) satisfies \(\varphi\) iff \(d'\) does.

(\(\Leftarrow\)) Let \(\Gamma\) be satisfiable, with any two data satisfying
\(\Gamma\) agreeing on every formula. Fix \(d_0\entails_{\mathcal S}\Gamma\)
and \(\varphi\in F\). If \(d_0\entails_{\mathcal S}\varphi\), every \(d\)
satisfying \(\Gamma\) satisfies \(\varphi\), so
\(\Gamma\entails_{\mathcal S}\varphi\). If not, every \(d\) satisfying
\(\Gamma\) fails \(\varphi\), and by the classical clause it satisfies
\(\neg\varphi\), so \(\Gamma\entails_{\mathcal S}\neg\varphi\). Hence
\(\Gamma\) is deductively complete. It is consistent by clause 1.

\emph{Clause 4.} Here the data are pairs \((\mathcal M,s)\) of an
\(\Lang\)-structure and an assignment. Consistency of \(T\) is
satisfiability by clause 1, and this is the existence of a model by clause 1
of \cref{thm:p01-w02-p01-016-sentences}. For a sentence \(\sigma\) and a
structure \(\mathcal M\), \(\mathcal M\entails\neg\sigma\) iff
\(\mathcal M\not\entails\sigma\). Indeed, by the assignment independence of
sentences (\cref{sec:fol-semantics-02}), \(\mathcal M,s\entails\sigma\) has
the same truth value for every assignment \(s\), and assignments exist
(clause 1 of \cref{thm:p01-w02-p01-016-variants}). So
\(\mathcal M\entails\neg\sigma\) iff some (equivalently every) \(s\) has
\(\mathcal M,s\not\entails\sigma\), iff \(\mathcal M\not\entails\sigma\). By
the displayed equivalence and clause 2 of
\cref{thm:p01-w02-p01-016-sentences}, the sentence condition reads: for every
sentence \(\sigma\), every model of \(T\) is a model of \(\sigma\), or every
model of \(T\) is a model of \(\neg\sigma\).

(\(\Rightarrow\)) Let \(\mathcal M,\mathcal N\) be models of \(T\) and
\(\sigma\) a sentence. In the first case both are models of \(\sigma\). In
the second case both are models of \(\neg\sigma\), so neither is a model of
\(\sigma\). Either way they agree on \(\sigma\).

(\(\Leftarrow\)) Let \(T\) have a model \(\mathcal M_0\), with all models of
\(T\) agreeing on every sentence. Let \(\sigma\) be a sentence. If
\(\mathcal M_0\entails\sigma\), every model of \(T\) is a model of
\(\sigma\). If not, no model of \(T\) is a model of \(\sigma\), so every model
of \(T\) is a model of \(\neg\sigma\). So the sentence condition holds, and
\(T\) is consistent because it has a model.
\end{proof}
\end{annexitem}

\begin{annexitem}{annex:p01-w02-p01-019-least-closed}{thm:p01-w02-p01-019-least-closed}{Derivability as least closed superset}
\begin{proof}
Write \(D(\Gamma)=\{\chi:\Gamma\derives_{\mathcal C_R}\chi\}\).

\emph{\(\Gamma\subseteq D(\Gamma)\).} For \(\chi\in\Gamma\), the one-node tree
whose node is an assumption leaf labelled \(\chi\) is a derivation with root
\(\chi\) and its only assumption leaf in \(\Gamma\). This is also
reflexivity.

\emph{\(D(\Gamma)\) is \(R\)-closed.} Let \((\Pi,\chi)\in R\) with
\(\Pi\subseteq D(\Gamma)\). For each \(\pi\in\Pi\) choose a derivation
\(E_\pi\) with root \(\pi\) and assumption leaves in \(\Gamma\). The set
\(\Pi\) is finite, so finitely many are chosen. Form the tree with a new root,
a rule node for \((\Pi,\chi)\) labelled \(\chi\), whose children are the roots
of the \(E_\pi\), one for each \(\pi\in\Pi\). If \(\Pi=\varnothing\), the tree
is the single rule node. It is a finite tree, and each node satisfies the
condition of \cref{def:p01-w02-p01-019-rule-calculus}. The new node is a rule
node, and every other node keeps its kind and children. No rule discharges an
assumption, so its assumption leaves are those of the \(E_\pi\), all in
\(\Gamma\). Hence \(\chi\in D(\Gamma)\).

\emph{Minimality.} Let \(X\supseteq\Gamma\) be \(R\)-closed. We show, by
induction on the height of a derivation \(E\) whose assumption leaves are in
\(\Gamma\), that its root label lies in \(X\). If the root is an assumption
leaf, its label is in \(\Gamma\subseteq X\). If the root is a rule node for
\((\Pi,\chi)\), its children are the roots of subtrees of smaller height. Each
subtree is a derivation whose assumption leaves are among those of \(E\), so
they lie in \(\Gamma\). By the induction hypothesis the children's labels,
which are exactly the members of \(\Pi\), lie in \(X\). Closure gives
\(\chi\in X\). This includes the case \(\Pi=\varnothing\). Hence
\(D(\Gamma)\subseteq X\).
\end{proof}
\end{annexitem}

\begin{annexitem}{annex:p01-w02-p01-019-three-calculi}{thm:p01-w02-p01-019-three-calculi}{Derivability and the properties of the three test calculi}
\begin{proof}
All three calculi have reflexivity by
\cref{thm:p01-w02-p01-019-least-closed}. The valuation semantics has the
classical negation clause, and no valuation satisfies \(\bottom\)
(\cref{sec:prop-syntax-semantics-02}). We write \(v\) for valuations, and
\(v_{\mathbf T}\), \(v_{\mathbf F}\) for the valuations sending every
propositional variable to \(\mathbf T\), respectively to \(\mathbf F\).
Distinct formulas in the claims below are distinct strings. For instance,
\(P\lor Q\) differs from \(P\) because it begins with a parenthesis.

\emph{Clause 1, \(\mathcal C_0\).} With \(R=\varnothing\), every set is
\(R\)-closed, \(\Gamma\) in particular. By
\cref{thm:p01-w02-p01-019-least-closed},
\(\Gamma\subseteq D(\Gamma)\subseteq\Gamma\), so
\(\Gamma\derives_{\mathcal C_0}\chi\) iff \(\chi\in\Gamma\).
\begin{itemize}
  \item \emph{Sound.} If \(\chi\in\Gamma\), every valuation satisfying
        \(\Gamma\) satisfies \(\chi\).
  \item \emph{Not explosive.} \(\{P,\neg P\}\) derives \(P\) and \(\neg P\),
        but not \(Q\notin\{P,\neg P\}\). This shows at once that
        \(\{P,\neg P\}\) is \(\mathcal C_0\)-inconsistent and absolutely
        \(\mathcal C_0\)-consistent.
  \item \emph{Unsatisfiable.} \(\{P,\neg P\}\) is unsatisfiable by the
        classical clause.
  \item \emph{Not complete.} \(\{P\}\entails P\lor Q\) by the disjunction
        clause, while \(P\lor Q\notin\{P\}\).
\end{itemize}

\emph{Clause 2, \(\mathcal C_{\mathrm{ex}}\).}
\begin{itemize}
  \item \emph{Right to left.} If \(\chi\in\Gamma\), use reflexivity. If
        \(\Gamma\) is unsatisfiable, clause 2 of
        \cref{thm:prop-compactness} gives a finite unsatisfiable
        \(\Gamma_0\subseteq\Gamma\). Then
        \((\Gamma_0,\chi)\in R_{\mathrm{ex}}\), and the tree with root a rule
        node for \((\Gamma_0,\chi)\) and one assumption leaf for each member
        of \(\Gamma_0\) derives \(\chi\) from \(\Gamma\).
  \item \emph{Left to right.} If \(\Gamma\) is unsatisfiable there is
        nothing to prove. If some \(v\) satisfies \(\Gamma\), then \(v\)
        satisfies every finite \(\Pi\subseteq\Gamma\). So no
        \((\Pi,\chi)\in R_{\mathrm{ex}}\) has \(\Pi\subseteq\Gamma\), and
        \(\Gamma\) is vacuously \(R_{\mathrm{ex}}\)-closed. By
        \cref{thm:p01-w02-p01-019-least-closed},
        \(\Gamma\derives_{\mathcal C_{\mathrm{ex}}}\chi\) implies
        \(\chi\in\Gamma\).
\end{itemize}
The properties follow from this characterization.
\begin{itemize}
  \item \emph{Sound.} If \(\chi\in\Gamma\) then \(\Gamma\entails\chi\). If
        \(\Gamma\) is unsatisfiable, then \(\Gamma\entails\chi\) vacuously
        (clause 4 of \cref{thm:p01-w02-p01-016-unsat-reduction}).
  \item \emph{(B1).} If \(\Gamma\derives\bottom\), then \(\bottom\in\Gamma\)
        or \(\Gamma\) is unsatisfiable. In the first case \(\Gamma\) is
        unsatisfiable too, since no valuation satisfies \(\bottom\). So
        \(\Gamma\) derives every formula.
  \item \emph{(B2).} If \(\Gamma\derives\varphi\) and
        \(\Gamma\derives\neg\varphi\), then either \(\Gamma\) is
        unsatisfiable or \(\varphi,\neg\varphi\in\Gamma\). In the second
        case \(\Gamma\) is unsatisfiable by the classical clause. Either
        way \(\Gamma\derives\bottom\).
  \item \emph{Explosion.} It follows by clause 4 of
        \cref{thm:p01-w02-p01-019-forms}.
  \item \emph{Consistent sets.} A satisfiable set is consistent by clause 1
        of \cref{thm:p01-w02-p01-019-soundness}. An unsatisfiable set
        derives \(P\) and \(\neg P\), so it is inconsistent. Hence the
        consistent sets are exactly the satisfiable sets.
  \item \emph{Not complete.} \(\{P\}\) is satisfiable (by
        \(v_{\mathbf T}\)) and \(P\lor Q\notin\{P\}\), so
        \(\{P\}\notderives_{\mathcal C_{\mathrm{ex}}}P\lor Q\), although
        \(\{P\}\entails P\lor Q\).
  \item \emph{No reductio.} The set \(\{P,\neg(P\lor Q)\}\) is
        unsatisfiable: a valuation satisfying \(P\) satisfies \(P\lor Q\),
        hence fails \(\neg(P\lor Q)\). So the set is inconsistent, but
        \(\{P\}\notderives_{\mathcal C_{\mathrm{ex}}}P\lor Q\). Reductio
        fails at \(\Gamma=\{P\}\), \(\varphi=P\lor Q\).
\end{itemize}

\emph{Clause 3, \(\mathcal C_P\).} Let
\(X=\{\chi:\Gamma\cup\{P\}\entails\chi\}\).
\begin{itemize}
  \item \emph{Right to left.} If \(\Gamma\cup\{P\}\entails\chi\), clause 1
        of \cref{thm:prop-compactness} gives a finite
        \(\Gamma_0\subseteq\Gamma\cup\{P\}\) with \(\Gamma_0\entails\chi\).
        Put \(\Pi=\Gamma_0\setminus\{P\}\). It is a finite subset of
        \(\Gamma\), and \(\Gamma_0\subseteq\Pi\cup\{P\}\), so
        \(\Pi\cup\{P\}\entails\chi\) by clause 3 of
        \cref{thm:p01-w02-p01-016-unsat-reduction}. Thus
        \((\Pi,\chi)\in R_P\), and the tree with root a rule node for
        \((\Pi,\chi)\) and one assumption leaf for each member of \(\Pi\)
        derives \(\chi\) from \(\Gamma\).
  \item \emph{Left to right.} \(\Gamma\subseteq X\). Also \(X\) is
        \(R_P\)-closed: if \(\Pi\subseteq X\) and
        \(\Pi\cup\{P\}\entails\chi\), let \(v\) satisfy \(\Gamma\cup\{P\}\).
        Then \(v\) satisfies each member of \(\Pi\), since each lies in
        \(X\), and it satisfies \(P\), so it satisfies \(\chi\). By
        \cref{thm:p01-w02-p01-019-least-closed}, derivable formulas lie in
        \(X\).
\end{itemize}
The properties follow from this characterization.
\begin{itemize}
  \item \emph{Complete.} If \(\Gamma\entails\chi\), then
        \(\Gamma\cup\{P\}\entails\chi\) by clause 3 of
        \cref{thm:p01-w02-p01-016-unsat-reduction}, so
        \(\Gamma\derives_{\mathcal C_P}\chi\).
  \item \emph{(B1).} If \(\Gamma\cup\{P\}\entails\bottom\), no valuation
        satisfies \(\Gamma\cup\{P\}\). Then \(\Gamma\cup\{P\}\entails\psi\)
        for every \(\psi\).
  \item \emph{(B2).} If \(\Gamma\cup\{P\}\) entails \(\varphi\) and
        \(\neg\varphi\), a satisfying valuation would satisfy both,
        contrary to the classical clause. So \(\Gamma\cup\{P\}\) is
        unsatisfiable and entails \(\bottom\).
  \item \emph{Explosion.} It follows by clause 4 of
        \cref{thm:p01-w02-p01-019-forms}.
  \item \emph{Reductio.} If \(\Gamma\cup\{\neg\varphi\}\) derives \(\psi\)
        and \(\neg\psi\), then \(\Gamma\cup\{\neg\varphi,P\}\) entails both.
        As in (B2) it is unsatisfiable. By clause 1 of
        \cref{thm:p01-w02-p01-016-unsat-reduction}, applied to the set
        \(\Gamma\cup\{P\}\), we get \(\Gamma\cup\{P\}\entails\varphi\), that
        is, \(\Gamma\derives_{\mathcal C_P}\varphi\).
  \item \emph{Consistent.} If \(\varnothing\) derived \(\varphi\) and
        \(\neg\varphi\), then \(\{P\}\) would entail both, and
        \(v_{\mathbf T}\), which satisfies \(P\), would satisfy both. This
        is impossible.
  \item \emph{Consistent sets are satisfiable.} This holds by clause 1 of
        \cref{thm:p01-w02-p01-019-completeness}.
  \item \emph{Not sound.} \(\{P\}\entails P\), so
        \(\derives_{\mathcal C_P}P\). But \(v_{\mathbf F}\) fails \(P\), so
        \(\not\entails P\).
  \item \emph{\(\{\neg P\}\).} It is satisfiable, by \(v_{\mathbf F}\). The
        set \(\{\neg P,P\}\) is unsatisfiable, so it entails every formula.
        In particular \(\{\neg P\}\) derives \(P\) and \(\neg P\) in
        \(\mathcal C_P\), and it is \(\mathcal C_P\)-inconsistent.
\end{itemize}
\end{proof}
\end{annexitem}

\begin{annexitem}{annex:p01-w02-p01-019-separation}{thm:p01-w02-p01-019-separation}{The three notions are independent}
\begin{proof}
\emph{Clauses 1--3.} These are read off from
\cref{thm:p01-w02-p01-019-three-calculi}.
\begin{itemize}
  \item \(\mathcal C_P\) is consistent, complete, has reductio and is not
        sound. The satisfiable set \(\{\neg P\}\) is
        \(\mathcal C_P\)-inconsistent, so the condition in the converse part
        of \cref{thm:p01-w02-p01-019-soundness} fails for it while it holds
        at \(\varnothing\).
  \item \(\mathcal C_{\mathrm{ex}}\) is sound and explosive, and its
        consistent sets are satisfiable, but it is not complete.
        \(\mathcal C_P\) is complete and not sound.
  \item \(\{P,\neg P\}\) is absolutely \(\mathcal C_0\)-consistent and
        \(\mathcal C_0\)-inconsistent, so the inclusion of clause 1 of
        \cref{thm:p01-w02-p01-019-forms} is strict for \(\mathcal C_0\).
        \(\mathcal C_0\) is sound and not complete.
\end{itemize}

\emph{Clause 4.} \(\mathsf{ND}_{\mathsf{CPL}}\) is sound
(\cref{thm:prop-nd-soundness}) and complete (\cref{thm:prop-completeness}).
By clause 2 of \cref{thm:p01-w02-p01-019-soundness}, since valuations exist,
\(\varnothing\) is consistent. If \(\derives P\), soundness would make \(P\)
valid, but \(v_{\mathbf F}\) fails \(P\). If \(\derives\neg P\), then
\(\neg P\) would be valid, but \(v_{\mathbf T}\) satisfies \(P\) and so fails
\(\neg P\). So \(\varnothing\) is not deductively complete.

For \(\Gamma_v\), the valuation \(v\) satisfies every \(p\in\Gamma_v\) with
\(v\entails p\), and every \(\neg p\in\Gamma_v\), since then \(v\not\entails p\).
So \(\Gamma_v\) is satisfiable. Let \(w\) satisfy \(\Gamma_v\), and let
\(p\in\mathsf{Prop}\). If \(v\entails p\), then \(p\in\Gamma_v\) and
\(w\entails p\). If \(v\not\entails p\), then \(\neg p\in\Gamma_v\) and
\(w\not\entails p\). So \(w\) and \(v\) agree on every propositional
variable. By \cref{lem:prop-coincidence} they agree on every formula. Any two
valuations satisfying \(\Gamma_v\) therefore agree with \(v\), and hence with
each other, on every formula. Clause 3 of
\cref{thm:p01-w02-p01-019-adequate}, applicable because the valuation
semantics has the classical clause, shows that \(\Gamma_v\) is consistent and
deductively complete for \(\mathsf{ND}_{\mathsf{CPL}}\).
\end{proof}
\end{annexitem}

\begin{annexitem}{annex:p01-w02-p01-019-intuitionistic}{thm:p01-w02-p01-019-intuitionistic}{Failure of the classical reduction for intuitionistic forcing}
\begin{proof}
The inherited definition says that \(\Gamma\entails_{\mathsf{IPL}}A\) holds iff,
for every structure \(\mathcal M\) and every world \(w\), forcing all members
of \(\Gamma\) at \(w\) implies forcing \(A\) at \(w\). This is the definition
of consequence in \cref{sec:levels-03} with the pairs \((\mathcal M,w)\) as
data. Since \(\neg A\) is \(A\to\bottom\) and no world forces \(\bottom\), the
implication clause gives
\[
  \mathcal M,w\forces\neg A
  \quad\Longleftrightarrow\quad
  \text{no } u\geq w \text{ has } \mathcal M,u\forces A .
\]

\emph{Clause 1.} If \(w\forces A\) and \(w\forces\neg A\), take \(u=w\geq w\)
in the display. This yields a contradiction.

\emph{Clause 2.} Suppose \(w\forces\neg(p\lor\neg p)\). By the display, no
\(u\geq w\) forces \(p\lor\neg p\). In particular, at \(u=w\), neither \(p\)
nor \(\neg p\) is forced. Since \(w\nforces\neg p\), the display gives some
\(u\geq w\) with \(u\forces p\). By the disjunction clause
\(u\forces p\lor\neg p\), contradicting the choice of \(w\). So no world
forces \(\neg(p\lor\neg p)\), and \(\{\neg(p\lor\neg p)\}\) is unsatisfiable.

\emph{Clause 3.} Let \(W=\{w_0,w_1\}\) with \(w_0<w_1\), the order being
reflexive, and let \(V(p)=\{w_1\}\). The value of \(V\) at other variables is
\(\varnothing\). Each value is upward closed, so \(V\) is persistent. Then
\(w_0\nforces p\). Also \(w_0\nforces\neg p\) by the display, because
\(w_1\geq w_0\) forces \(p\). By the disjunction clause,
\(w_0\nforces p\lor\neg p\). The world \(w_0\) forces every member of
\(\varnothing\), so \(\varnothing\not\entails_{\mathsf{IPL}}p\lor\neg p\).
This is the structure of the inherited counterexample
\cite[ch.~``Intuitionistic Logic'', example ``Counterexample structure and
world'']{FrozenFormalLogicConsistencyComparison}.

\emph{Clause 4.} Let \(\mathcal C\) be sound and complete for this semantics.
By clause 2 the set \(\{\neg(p\lor\neg p)\}\) is unsatisfiable. Clause 1 of
\cref{thm:p01-w02-p01-019-completeness} holds for every semantics and so
makes it trivial. It therefore derives \(p\lor\neg p\) and
\(\neg(p\lor\neg p)\), and it is \(\mathcal C\)-inconsistent. If
\(\varnothing\derives_{\mathcal C}p\lor\neg p\), soundness would give
\(\varnothing\entails_{\mathsf{IPL}}p\lor\neg p\), contrary to clause 3.
Reductio with \(\Gamma=\varnothing\) and \(\varphi=p\lor\neg p\) would
require that derivation, so \(\mathcal C\) lacks reductio.
\end{proof}
\end{annexitem}
